\section{Especificaciones de casos de uso}

\subsection{Controles de simulación}

Los siguientes casos de uso permiten al usuario controlar la evolución del autómata celular.

\subsubsection{Iniciar simulación}

\begin{table}[H]
\begin{center}
	\begin{tabular}{| p{0.2\textwidth} | p{0.7\textwidth} |}
		\hline 
			\textbf{Actor} & Usuario\\ \hline
			\textbf{Descripción} & El sistema inicia la evolución en pasos de tiempo constante. \\ \hline
			\textbf{Tipo} & Primario \\ \hline
			\textbf{Entradas} &
			\begin{itemize}[noitemsep, nosep]
			\item Desde mouse \BTN{iniciar_simulacion}
				\end{itemize}
			\\ \hline
			\textbf{Salidas} & Ninguna.  \\ \hline
			\textbf{Precondición} & \begin{itemize}[noitemsep, nosep]
				\item Simulación detenida.
				\item Botón \BTN{iniciar_simulacion} habilitado.
			\end{itemize}
			\\ \hline
			\textbf{Postcondición} & \begin{itemize}[noitemsep, nosep]
				\item Simulación en ejecución.
				\item Botón \BTN{iniciar_simulacion} deshabilitado.
			\end{itemize} \\ \hline
			\textbf{Errores}  & Ninguno.  \\ \hline
			\end{tabular}
		\end{center}
	\end{table}

\begin{table}[H]
\begin{center}
	\begin{tabular}{|p{0.45\textwidth}|p{0.45\textwidth}|}
		\hline
			\multicolumn{2}{|c|}{\textbf{Flujo básico}} \\ \hline
			\textbf{Actor} & \textbf{Sistema} \\
			\hline
			\begin{enumerate}[noitemsep, nosep]
			\item Presiona el botón \BTN{iniciar_simulacion} desde \UI{control_ui}.
				\end{enumerate} 
			& 
			\begin{enumerate}[noitemsep, nosep]
			\item[2.] Inicia la simulación. Cada paso evolutivo se realiza a una velocidad constante. 
				\item Inhabilita el botón \BTN{iniciar_simulacion}.
				\item Habilita el botón \BTN{pausar_simulacion}.
				\end{enumerate} \\
			\hline
			\end{tabular} \\
		\end{center}
	\end{table}

\subsubsection{Pausar simulación}

\begin{table}[H]
\begin{center}
	\begin{tabular}{| p{0.2\textwidth} | p{0.7\textwidth} |}
		\hline 
			\textbf{Actor} & Usuario\\ \hline
			\textbf{Descripción} & Se detiene la simulación.  \\ \hline
			\textbf{Tipo} & Secundario \\ \hline
			\textbf{Entradas} &
			\begin{itemize}[noitemsep, nosep]
			\item Desde mouse \BTN{pausar_simulacion}.
				\end{itemize}
			\\ \hline
			\textbf{Salidas} & Ninguna. \\ \hline
			\textbf{Precondición} &  \begin{itemize}[noitemsep, nosep]
				\item Simulación en ejecución.
				\item Botón \BTN{pausar_simulacion} habilitado.
			\end{itemize}			\\ \hline
			\textbf{Postcondición} & \begin{itemize}[noitemsep, nosep]
				\item Simulación detenida.
				\item Botón \BTN{pausar_simulacion} deshabilitado.
			\end{itemize}  \\ \hline
			\textbf{Errores} & Ninguno. \\ \hline
			\end{tabular}
			\end{center}
				\end{table}
			
		\begin{table}[H]
	\begin{center}
\begin{tabular}{|p{0.45\textwidth}|p{0.45\textwidth}|}
\hline
\multicolumn{2}{|c|}{\textbf{Flujo básico}} \\ \hline
	\textbf{Actor} & \textbf{Sistema} \\
		\hline
			\begin{enumerate}[noitemsep, nosep]
			\item Presiona el botón \BTN{pausar_simulacion} desde \UI{control_ui}.
			\end{enumerate} 
			& 
			\begin{enumerate}[noitemsep, nosep]
				\item[2.] Detiene la simulación. 
			\item[3.] Inhabilita el botón \BTN{pausar_simulacion}.
			\item[4.] Habilita el botón \BTN{iniciar_simulacion}. 
			\end{enumerate} \\
				\hline
			\end{tabular} \\
			\end{center}
		\end{table}
	
\subsubsection{Aumentar velocidad}

\begin{table}[H]
	\begin{center}
		\begin{tabular}{| p{0.2\textwidth} | p{0.7\textwidth} |}
			\hline 
			\textbf{Actor} & Usuario\\ \hline
			\textbf{Descripción} & Se aumenta la velocidad de actualización del autómata celular. \\ \hline
			\textbf{Tipo} & Primario. \\ \hline
			\textbf{Entradas} &
			\begin{itemize}[noitemsep, nosep]
				\item Desde mouse \BTN{aumentar_velocidad}
			\end{itemize}
			\\ \hline
			\textbf{Salidas} & Ninguna \\ \hline
			\textbf{Precondición} & Ninguna.
			\\ \hline
			\textbf{Postcondición} & \begin{itemize}[noitemsep, nosep]
				\item Aumenta la velocidad de simulación.
			\end{itemize}  \\ \hline
			\textbf{Errores} & 
			\begin{itemize}[noitemsep, nosep]
				\item Velocidad límite: Se intenta aumentar la velocidad cuando ya encuentra en su velocidad máxima.
			\end{itemize}\\ \hline
		\end{tabular}
	\end{center}
\end{table}

\begin{table}[H]
	\begin{center}
		\begin{tabular}{|p{0.45\textwidth}|p{0.45\textwidth}|}
			\hline
			\multicolumn{2}{|c|}{\textbf{Flujo básico}} \\ \hline
			\textbf{Actor} & \textbf{Sistema} \\
			\hline 
			\begin{enumerate}[noitemsep, nosep]
				\item Presiona \BTN{aumentar_velocidad} desde \UI{control_ui}.
			\end{enumerate} 
			& 
			\begin{enumerate}[noitemsep, nosep] 
				\item[2.] Aumenta la velocidad de la simulación.
			\end{enumerate} \\
			\hline
		\end{tabular} \\
	\end{center}
\end{table}

\begin{table}[H]
	\begin{center}
		\begin{tabular}{|p{0.45\textwidth}|p{0.45\textwidth}|}
			\hline
			\multicolumn{2}{|c|}{\textbf{Flujo alterno 1: Velocidad límite}} \\ \hline
			\textbf{Actor} & \textbf{Sistema} \\
			\hline
			\begin{enumerate}[noitemsep, nosep]
				\item[1.1.] Presiona \BTN{aumentar_velocidad} desde \UI{control_ui} cuando está en la velocidad límite máxima. 
			\end{enumerate}
			& 
			\begin{enumerate}[noitemsep, nosep]
				\item[1.2.] No realiza ningún cambio, mantiene la velocidad actual. 
			\end{enumerate} \\
			\hline
		\end{tabular} \\
	\end{center}
\end{table}


\subsubsection{Reducir velocidad}

\begin{table}[H]
	\begin{center}
		\begin{tabular}{| p{0.2\textwidth} | p{0.7\textwidth} |}
			\hline 
			\textbf{Actor} & Usuario\\ \hline
			\textbf{Descripción} & Se disminuye la velocidad de actualización del autómata celular. \\ \hline
			\textbf{Tipo} & Primario. \\ \hline
			\textbf{Entradas} &
			\begin{itemize}[noitemsep, nosep]
				\item Desde mouse \BTN{disminuir_velocidad}
			\end{itemize}
			\\ \hline
			\textbf{Salidas} & Ninguna \\ \hline
			\textbf{Precondición} & Ninguna.
			\\ \hline
			\textbf{Postcondición} & \begin{itemize}[noitemsep, nosep]
				\item Disminuye la velocidad de simulación.
			\end{itemize}  \\ \hline
			\textbf{Errores} & 
			\begin{itemize}[noitemsep, nosep]
				\item Velocidad límite: Se intenta disminuir la velocidad cuando ya encuentra en su velocidad mínima.
			\end{itemize}\\ \hline
		\end{tabular}
	\end{center}
\end{table}

\begin{table}[H]
	\begin{center}
		\begin{tabular}{|p{0.45\textwidth}|p{0.45\textwidth}|}
			\hline
			\multicolumn{2}{|c|}{\textbf{Flujo básico}} \\ \hline
			\textbf{Actor} & \textbf{Sistema} \\
			\hline 
			\begin{enumerate}[noitemsep, nosep]
				\item Presiona \BTN{disminuir_velocidad} desde \UI{control_ui}.
			\end{enumerate} 
			& 
			\begin{enumerate}[noitemsep, nosep] 
				\item[2.] Disminuye la velocidad de la simulación.
			\end{enumerate} \\
			\hline
		\end{tabular} \\
	\end{center}
\end{table}

\begin{table}[H]
	\begin{center}
		\begin{tabular}{|p{0.45\textwidth}|p{0.45\textwidth}|}
			\hline
			\multicolumn{2}{|c|}{\textbf{Flujo alterno 1: Velocidad límite}} \\ \hline
			\textbf{Actor} & \textbf{Sistema} \\
			\hline
			\begin{enumerate}[noitemsep, nosep]
				\item[1.1.] Presiona \BTN{disminuir_velocidad} desde \UI{control_ui} cuando está en la velocidad límite mínima. 
			\end{enumerate}
			& 
			\begin{enumerate}[noitemsep, nosep]
				\item[1.2.] No realiza ningún cambio, mantiene la velocidad actual. 
			\end{enumerate} \\
			\hline
		\end{tabular} \\
	\end{center}
\end{table}

\subsubsection{Simular un paso}

\begin{table}[H]
\begin{center}
\begin{tabular}{| p{0.2\textwidth} | p{0.7\textwidth} |}
	\hline 
		\textbf{Actor} & Usuario\\ \hline
			\textbf{Descripción} & Se realiza la simulación del CA, pero únicamente en un paso. \\ \hline
			\textbf{Tipo} & Primario \\ \hline
			\textbf{Entradas} &
			\begin{itemize}[noitemsep, nosep]
			\item Desde mouse \BTN{paso_adelante}
			\end{itemize}
				\\ \hline
			\textbf{Salidas} & Ninguna. \\ \hline
			\textbf{Precondición} & Ninguna.
			\\ \hline
			\textbf{Postcondición} & \begin{itemize}[noitemsep, nosep]
				\item Evoluciona el sistema.
			\end{itemize}  \\ \hline
			\textbf{Errores} & 
			\begin{itemize}[noitemsep, nosep]
			\item Velocidad máxima: Se intenta dar un paso cuando la simulación está activa en la velocidad máxima.
			\end{itemize}\\ \hline
				\end{tabular}
			\end{center}
		\end{table}
	
\begin{table}[H]
\begin{center}
\begin{tabular}{|p{0.45\textwidth}|p{0.45\textwidth}|}
	\hline
		\multicolumn{2}{|c|}{\textbf{Flujo básico}} \\ \hline
			\textbf{Actor} & \textbf{Sistema} \\
			\hline
			\begin{enumerate}[noitemsep, nosep]
			\item Presiona el botón \BTN{paso_adelante} desde \UI{control_ui}.
			\end{enumerate} 
				& 
			\begin{enumerate}[noitemsep, nosep]
			\item[2.] La simulación avanza un paso.
			\end{enumerate} \\
				\hline
			\end{tabular} \\
			\end{center}
		\end{table}
	
\begin{table}[H]
\begin{center}
\begin{tabular}{|p{0.45\textwidth}|p{0.45\textwidth}|}
	\hline
		\multicolumn{2}{|c|}{\textbf{Flujo alterno 1: Velocidad máxima}} \\ \hline
			\textbf{Actor} & \textbf{Sistema} \\
			\hline
			\begin{enumerate}[noitemsep, nosep]
			\item[1.1.] Presiona el botón \BTN{paso_adelante} desde \UI{control_ui} cuando la velocidad esta en la velocidad máxima.
			\end{enumerate} 
				& 
			\begin{enumerate}[noitemsep, nosep]
			\item[1.2.] No realiza ningún cambio y mantiene la velocidad actual. 
			\end{enumerate} \\
				\hline
			\end{tabular} 
			\end{center}
		\end{table}
	
\subsubsection{Retroceder un paso}


\begin{table}[H]
	\begin{center}
		\begin{tabular}{| p{0.2\textwidth} | p{0.7\textwidth} |}
			\hline 
			\textbf{Actor} & Usuario\\ \hline
			\textbf{Descripción} & La simulación adquiere el estado anterior al actual.  \\ \hline
			\textbf{Tipo} & Primario. \\ \hline
			\textbf{Entradas} &
			\begin{itemize}[noitemsep, nosep]
				\item Desde mouse \BTN{paso_atras}
			\end{itemize}
			\\ \hline
			\textbf{Salidas} & Ninguna. \\ \hline
			\textbf{Precondición} & Ninguna.
			\\ \hline
			\textbf{Postcondición} & \begin{itemize}[noitemsep, nosep]
				\item El sistema adquiere la evolución pasada.
			\end{itemize}  \\ \hline
			\textbf{Errores} & 
			\begin{itemize}[noitemsep, nosep]
				\item Configuración inicial: EL usuario intenta retroceder la simulación de la configuración inicial.
			\end{itemize}\\ \hline
		\end{tabular}
	\end{center}
\end{table}

\begin{table}[H]
	\begin{center}
		\begin{tabular}{|p{0.45\textwidth}|p{0.45\textwidth}|}
			\hline
			\multicolumn{2}{|c|}{\textbf{Flujo básico}} \\ \hline
			\textbf{Actor} & \textbf{Sistema} \\
			\hline
			\begin{enumerate}[noitemsep, nosep]
				\item Presiona el botón \BTN{paso_atras} desde \UI{control_ui}.
			\end{enumerate} 
			& 
			\begin{enumerate}[noitemsep, nosep]
				\item[2.] La simulación retrocede un paso. 
			\end{enumerate} \\
			\hline
		\end{tabular} \\
	\end{center}
\end{table}

\begin{table}[H]
	\begin{center}
		\begin{tabular}{|p{0.45\textwidth}|p{0.45\textwidth}|}
			\hline
			\multicolumn{2}{|c|}{\textbf{Flujo alterno 1: Configuración inicial}} \\ \hline
			\textbf{Actor} & \textbf{Sistema} \\
			\hline
			\begin{enumerate}[noitemsep, nosep]
				\item[1.1.] Presiona el botón \BTN{paso_atras} desde \UI{control_ui} estando en la configuración inicial (no hay más configuraciones antes).  
			\end{enumerate}
			& 
			\begin{enumerate}[noitemsep, nosep]
				\item[1.2.] No realiza ningún cambio.
			\end{enumerate} \\
			\hline
		\end{tabular} 
	\end{center}
\end{table}

\subsubsection{Reiniciar simulación}

\begin{table}[H]
	\begin{center}
		\begin{tabular}{| p{0.2\textwidth} | p{0.7\textwidth} |}
			\hline 
			\textbf{Actor} & Usuario\\ \hline
			\textbf{Descripción} & Se reinicia la simulación a la última cuadrícula que fue modificada mientras la simulación estaba en pausa. \\ \hline
			\textbf{Tipo} & Primario \\ \hline
			\textbf{Entradas} &
			\begin{itemize}[noitemsep, nosep]
				\item Desde mouse \BTN{reiniciar_simulacion}.
			\end{itemize}
			\\ \hline
			\textbf{Salidas} & Ninguna \\ \hline
			\textbf{Precondición} & Ninguna
			\\ \hline
			\textbf{Postcondición} & \begin{itemize}[noitemsep, nosep]
				\item El sistema adquiere la configuración inicial.
				\item El sistema adquiere la configuración nula.
			\end{itemize}  \\ \hline
			\textbf{Errores} & 
			\begin{itemize}[noitemsep, nosep]
				\item Configuración inicia: El usuario intenta reiniciar la simulación de la configuración inicial.
			\end{itemize}\\ \hline
		\end{tabular}
	\end{center}
\end{table}

\begin{table}[H]
	\begin{center}
		\begin{tabular}{|p{0.45\textwidth}|p{0.45\textwidth}|}
			\hline
			\multicolumn{2}{|c|}{\textbf{Flujo básico}} \\ \hline
			\textbf{Actor} & \textbf{Sistema} \\
			\hline
			\begin{enumerate}[noitemsep, nosep]
				\item Presiona el botón \BTN{reiniciar_simulacion} desde \UI{control_ui}.
			\end{enumerate} 
			& 
			\begin{enumerate}[noitemsep, nosep]
				\item[2.] Regresa la configuración a la última versión modificada mientras estaba en pausa. 
			\end{enumerate} \\
			\hline
		\end{tabular} \\
	\end{center}
\end{table}

\begin{table}[H]
	\begin{center}
		\begin{tabular}{|p{0.45\textwidth}|p{0.45\textwidth}|}
			\hline
			\multicolumn{2}{|c|}{\textbf{Flujo alterno 1: Configuración inicial}} \\ \hline
			\textbf{Actor} & \textbf{Sistema} \\
			\hline
			\begin{enumerate}[noitemsep, nosep]
				\item[1.1.] Presiona el botón \BTN{reiniciar_simulacion} desde \UI{control_ui} estando en la configuración inicial (no hay más configuraciones antes).  
			\end{enumerate}
			& 
			\begin{enumerate}[noitemsep, nosep]
				\item[1.2.] Establece todas las células a 0.
			\end{enumerate} \\
			\hline
		\end{tabular} 
	\end{center}
\end{table}


\subsection{Configuración del autómata celular}

Estos casos de uso definen las características del autómata celular, como el numero de estados, vecindad y la función de transición.

\subsubsection{Definir vecindad}

\begin{table}[H]
\begin{center}
	\begin{tabular}{| p{0.2\textwidth} | p{0.7\textwidth} |}
		\hline 
			\textbf{Actor} & Usuario \\ \hline
			\textbf{Descripción} & Se define el conjunto de posiciones relativas de cada célula que forman parte de la vecindad.  \\ \hline
			\textbf{Tipo} & Primario \\ \hline
			\textbf{Entradas} &
			\begin{itemize}[noitemsep, nosep]
			\item Desde mouse \BTN{def_vecindad}
			\item Desde mouse \BTN{vecindad_guardar}
			\item Desde mouse \BTN{vecindad_cancelar}
			\end{itemize}
			\\ \hline
			\textbf{Salidas} & 
			\begin{itemize}[noitemsep, nosep]
			\item Mensaje \MSG{msg_vecindad}.
			\item Mensaje \MSG{msg_vecindad_incompatibilidad}
		\end{itemize} \\ \hline
		\textbf{Precondición} & Ninguna
		\\ \hline
		\textbf{Postcondición} &  \begin{itemize}[noitemsep, nosep]
			\item Vecindad definida.
		\end{itemize} \\ \hline
		\textbf{Errores} & 
	\begin{itemize}[noitemsep, nosep]
			\item Simulando: Se intenta cambiar la vecindad mientras está una simulación en proceso.
			\item Cancelar: Se cancela el cambio de vecindad.
			\item Incompatibilidad: La regla es incompatible con la nueva vecindad seleccionada.
			\end{itemize}\\ \hline
			
			\end{tabular}
				\end{center}
			\end{table}
			
\begin{table}[H]
\begin{center}
\begin{tabular}{|p{0.45\textwidth}|p{0.45\textwidth}|}
\hline
\multicolumn{2}{|c|}{\textbf{Flujo básico}} \\ \hline
\textbf{Actor} & \textbf{Sistema} \\
\hline
\begin{enumerate}[noitemsep, nosep]
\item Presiona el botón \BTN{def_vecindad} en \UI{configuracion_ui}.
\item[3. ] Selecciona los elementos que forman parte de la vecindad seleccionando las células.
\item[4. ] Presiona el botón \BTN{vecindad_guardar} en \UI{vecindad_ui}.
\end{enumerate} 
& 
\begin{enumerate}[noitemsep, nosep]
\item[2.] Abre la interfaz \UI{vecindad_ui}. 
\item[5.] Actualiza la vecindad.
\end{enumerate} \\
\hline
\end{tabular} 
\end{center}
\end{table}

\begin{table}[H]
\begin{center}
\begin{tabular}{|p{0.45\textwidth}|p{0.45\textwidth}|}
\hline
\multicolumn{2}{|c|}{\textbf{Flujo alterno 1: Simulando}} \\ \hline
\textbf{Actor} & \textbf{Sistema} \\
\hline
\begin{enumerate}[noitemsep, nosep]
\item[1.1.] Presiona el botón \BTN{def_vecindad} en \UI{configuracion_ui} mientras la simulación está activa.
\item[1.4.]  Selecciona los elementos que forman parte de la vecindad seleccionando las células.
\item[1.5.] Presiona el botón \BTN{vecindad_guardar} en \UI{vecindad_ui}.
\end{enumerate}
& 
\begin{enumerate}[noitemsep, nosep]
\item[1.2.] Detiene la simulación.
\item[1.3.] Abre la interfaz \UI{vecindad_ui}.
\item[1.6.] Actualiza la vecindad.
\end{enumerate} \\
\hline
\end{tabular} 
\end{center}
\end{table}

\begin{table}[H]
\begin{center}
\begin{tabular}{|p{0.45\textwidth}|p{0.45\textwidth}|}
\hline
\multicolumn{2}{|c|}{\textbf{Flujo alterno 2: Cancelar}} \\ \hline
\textbf{Actor} & \textbf{Sistema} \\
\hline
\begin{enumerate}[noitemsep, nosep]
\item[2.4.] Presiona el botón de \BTN{vecindad_cancelar} en \UI{vecindad_ui}. 
\end{enumerate}
& 
\begin{enumerate}[noitemsep, nosep]
\item[2.5.] No se actualiza la vecindad. 
\end{enumerate} \\
\hline
\end{tabular} 
\end{center}
\end{table}

\begin{table}[H]
\begin{center}
\begin{tabular}{|p{0.45\textwidth}|p{0.45\textwidth}|}
\hline
\multicolumn{2}{|c|}{\textbf{Flujo alterno 3: Incompatibilidad (Aceptar)}} \\ \hline
\textbf{Actor} & \textbf{Sistema} \\
\hline
\begin{enumerate}[noitemsep, nosep]
\item[3.3.] Selecciona los elementos que forman parte de la vecindad presionando los botones \BTN{celula} pero el nuevo tamaño es diferente al anterior. 
\item[3.4.] Presiona el botón \BTN{vecindad_guardar} en \UI{vecindad_ui}.
\item[3.6.] Acepta el cambio
\end{enumerate}
& 
\begin{enumerate}[noitemsep, nosep]
\item[3.5.] Muestra el mensaje \MSG{msg_vecindad_incompatibilidad}.
\item[3.7.] Actualiza la vecindad.
\end{enumerate} \\
\hline
\end{tabular} 
\end{center}
\end{table}

\begin{table}[H]
	\begin{center}
		\begin{tabular}{|p{0.45\textwidth}|p{0.45\textwidth}|}
			\hline
			\multicolumn{2}{|c|}{\textbf{Flujo alterno 3: Incompatibilidad (Rechazar)}} \\ \hline
			\textbf{Actor} & \textbf{Sistema} \\
			\hline
			\begin{enumerate}[noitemsep, nosep]
				\item[3.3.] Selecciona los elementos que forman parte de la vecindad presionando los botones \BTN{celula} pero el nuevo tamaño es diferente al anterior. 
				\item[3.4.] Presiona el botón \BTN{vecindad_guardar}.
				\item[3.6.] Rechaza el cambio
			\end{enumerate}
			& 
			\begin{enumerate}[noitemsep, nosep]
				\item[3.5.] Muestra el mensaje \MSG{msg_vecindad_incompatibilidad}.
				\item[3.7.] No actualiza la vecindad.
			\end{enumerate} \\
			\hline
		\end{tabular} 
	\end{center}
\end{table}

\subsubsection{Agregar estado}

\begin{table}[H]
	\begin{center}
		\begin{tabular}{| p{0.2\textwidth} | p{0.7\textwidth} |}
			\hline 
			\textbf{Actor} & Usuario \\ \hline
			\textbf{Descripción} & Se añade una nueva vecindad semitotalistica.  \\ \hline
			\textbf{Tipo} & Secundario \\ \hline
			\textbf{Entradas} &
			\begin{itemize}[noitemsep, nosep]
				\item Desde mouse \BTN{def_vecindad}
				\item Desde mouse \BTN{vecindad_añadir}
			\end{itemize}
			\\ \hline
			\textbf{Salidas} & 
			\begin{itemize}[noitemsep, nosep]
				\item Mensaje \MSG{msg_vecindad}.
				\item Mensaje \MSG{msg_vecindad_incompatibilidad}
			\end{itemize} \\ \hline
			\textbf{Precondición} & Ninguna
			\\ \hline
			\textbf{Postcondición} &  \begin{itemize}[noitemsep, nosep]
				\item Nueva vecindad.
			\end{itemize} \\ \hline
			\textbf{Errores} & Ninguno. \\ \hline
			
		\end{tabular}
	\end{center}
\end{table}

\begin{table}[H]
	\begin{center}
		\begin{tabular}{|p{0.45\textwidth}|p{0.45\textwidth}|}
			\hline
			\multicolumn{2}{|c|}{\textbf{Flujo básico}} \\ \hline
			\textbf{Actor} & \textbf{Sistema} \\
			\hline
			\begin{enumerate}[noitemsep, nosep]
				\item Presiona el botón \BTN{def_vecindad} en \UI{configuracion_ui}.
				\item[3. ] Presiona el botón \BTN{vecindad_añadir} en \UI{vecindad_ui}
			\end{enumerate} 
			& 
			\begin{enumerate}[noitemsep, nosep]
				\item[2.] Abre la interfaz \UI{vecindad_ui}. 
				\item[4.] Añade una nueva vecindad semitotalistica.
			\end{enumerate} \\
			\hline
		\end{tabular} 
	\end{center}
\end{table}

\subsubsection{Definir condición de frontera}

\begin{table}[H]
	\begin{center}
		\begin{tabular}{| p{0.2\textwidth} | p{0.7\textwidth} |}
			\hline 
			\textbf{Actor} & Usuario\\ \hline
			\textbf{Descripción} & Se selecciona entre diferentes opciones de condiciones de frontera. \\ \hline
			\textbf{Tipo} & Primario \\ \hline
			\textbf{Entradas} &
			\begin{itemize}[noitemsep, nosep]
				\item Desde mouse \BTN{def_frontera}
			\end{itemize}
			\\ \hline
			\textbf{Salidas} & Ninguna
			\\ \hline
			\textbf{Precondición} & 
			\begin{itemize}[noitemsep, nosep]
				\item Cambio de condición de frontera.
			\end{itemize}
			\\ \hline
			\textbf{Postcondición} & Ninguna \\ \hline
			\textbf{Errores} &Ninguno \\ \hline
		\end{tabular}
	\end{center}
\end{table}

\begin{table}[H]
	\begin{center}
		\begin{tabular}{|p{0.45\textwidth}|p{0.45\textwidth}|}
			\hline
			\multicolumn{2}{|c|}{\textbf{Flujo básico}} \\ \hline
			\textbf{Actor} & \textbf{Sistema} \\
			\hline
			\begin{enumerate}[noitemsep, nosep]
				\item Presiona el botón \BTN{frontera} desde \UI{configuracion_ui}.
			\end{enumerate} 
			& 
			\begin{enumerate}[noitemsep, nosep]
				\item[2.] Actualiza la condición de frontera a la siguiente configuración.  
			\end{enumerate} \\
			\hline
		\end{tabular} \\
	\end{center}
\end{table}

\subsubsection{Definir numero de estados}

\begin{table}[H]
	\begin{center}
		\begin{tabular}{| p{0.2\textwidth} | p{0.7\textwidth} |}
			\hline 
			\textbf{Actor} & Usuario\\ \hline
			\textbf{Descripción} & Se define el tamaño del conjunto de estados. \\ \hline
			\textbf{Tipo} & Primario \\ \hline
			\textbf{Entradas} &
			\begin{itemize}[noitemsep, nosep]
				\item Desde teclado \IN{estados} 
				\item Desde mouse \BTN{cambiar_estados}
			\end{itemize}
			\\ \hline
			\textbf{Salidas} & Ninguna  \\ \hline
			\textbf{Precondición} & Ninguna  
			\\ \hline
			\textbf{Postcondición} & Ninguna   \\ \hline
			\textbf{Errores} & 
			\begin{itemize}[noitemsep, nosep]
				\item Simulando: El usuario intenta cambiar los estados mientras está en ejecución.
				\item Incompatibilidad: El número de estados es menor y hay estados inexistentes en el espacio celular.
			\end{itemize}\\ \hline
		\end{tabular}
	\end{center}
\end{table}

\begin{table}[H]
	\begin{center}
		\begin{tabular}{|p{0.45\textwidth}|p{0.45\textwidth}|}
			\hline
			\multicolumn{2}{|c|}{\textbf{Flujo básico}} \\ \hline
			\textbf{Actor} & \textbf{Sistema} \\
			\hline
			\begin{enumerate}[noitemsep, nosep]
				\item Ingresa el número de estados en el campo \IN{estados} desde \UI{configuracion_ui}.
				\item Presiona el botón \BTN{cambiar_estados}.
			\end{enumerate} 
			& 
			\begin{enumerate}[noitemsep, nosep]
				\item[3.] Cambia el conjunto de estados.  
			\end{enumerate} \\
			\hline
		\end{tabular} \\
	\end{center}
\end{table}

\begin{table}[H]
	\begin{center}
		\begin{tabular}{|p{0.45\textwidth}|p{0.45\textwidth}|}
			\hline
			\multicolumn{2}{|c|}{\textbf{Flujo alterno 1: Simulando}} \\ \hline
			\textbf{Actor} & \textbf{Sistema} \\
			\hline
			\begin{enumerate}[noitemsep, nosep]
				\item[2.1.] Presiona el botón \BTN{cambiar_estados} mientras la simulación está en ejecución. 
			\end{enumerate}
			& 
			\begin{enumerate}[noitemsep, nosep]
				\item[3.1.] Detiene la simulación.
				\item[4.1.] Cambia el conjunto de estados.
			\end{enumerate} \\
			\hline
		\end{tabular} 
	\end{center}
\end{table}

\begin{table}[H]
	\begin{center}
		\begin{tabular}{|p{0.45\textwidth}|p{0.45\textwidth}|}
			\hline
			\multicolumn{2}{|c|}{\textbf{Flujo alterno 2: Incompatibilidad (Aceptar)}} \\ \hline
			\textbf{Actor} & \textbf{Sistema} \\
			\hline
			\begin{enumerate}[noitemsep, nosep]
				\item[1.2.] Ingresa el número de estados en el campo \IN{estados} desde \UI{configuracion_ui}, menor al actual.
				\item[2.2.] Presiona el botón \BTN{cambiar_estados}.
				\item[4.2.] Presiona el botón \BTN{aceptar_estados}.
			\end{enumerate}
			& 
			\begin{enumerate}[noitemsep, nosep]
				\item[3.2.] Muestra el mensaje \MSG{incompatibilidad_estados}.
				\item[5.2.] Cambia el conjunto de estados.
				\item[6.2.] Establece a ``0'' todos los estados inexistentes.
			\end{enumerate} \\
			\hline
		\end{tabular} 
	\end{center}
\end{table}

\begin{table}[H]
	\begin{center}
		\begin{tabular}{|p{0.45\textwidth}|p{0.45\textwidth}|}
			\hline
			\multicolumn{2}{|c|}{\textbf{Flujo alterno 2.1: Incompatibilidad (Rechazar)}} \\ \hline
			\textbf{Actor} & \textbf{Sistema} \\
			\hline
			\begin{enumerate}[noitemsep, nosep]
				\item[1.2.] Ingresa el número de estados en el campo \IN{estados} desde \UI{configuracion_ui} menor al actual.
				\item[2.2.] Presiona el botón \BTN{cambiar_estados}.
				\item[4.2.] Presiona el botón \BTN{rechazar_estados}.
			\end{enumerate}
			& 
			\begin{enumerate}[noitemsep, nosep]
				\item[3.2.] Muestra el mensaje \MSG{incompatibilidad_estados}.
				\item[5.2.] No hace ningún cambio.
			\end{enumerate} \\
			\hline
		\end{tabular} 
	\end{center}
\end{table}

\subsubsection{Definir función de transición}


\begin{table}[H]
	\begin{center}
		\begin{tabular}{| p{0.2\textwidth} | p{0.7\textwidth} |}
			\hline 
			\textbf{Actor} & Usuario\\ \hline
			\textbf{Descripción} & Se configura la regla que define el comportamiento del CA. \\ \hline
			\textbf{Tipo} & Primario \\ \hline
			\textbf{Entradas} &
			\begin{itemize}[noitemsep, nosep]
				\item Desde teclado \IN{regla}
				\item Desde mouse \BTN{aplicar_regla}
			\end{itemize}
			\\ \hline
			\textbf{Salidas} & Ninguna  \\ \hline
			\textbf{Precondición} & Ninguna 
			\\ \hline
			\textbf{Postcondición} &  Ninguna  \\ \hline
			\textbf{Errores} & 
			\begin{itemize}[noitemsep, nosep]
				\item Simulando: El cambio de regla se intenta aplicar cuando la simulación está en ejecución.
				\item Fuera de rango: La regla establecida está fuera del rango. 
			\end{itemize}\\ \hline
		\end{tabular}
	\end{center}
\end{table}

\begin{table}[H]
	\begin{center}
		\begin{tabular}{|p{0.45\textwidth}|p{0.45\textwidth}|}
			\hline
			\multicolumn{2}{|c|}{\textbf{Flujo básico}} \\ \hline
			\textbf{Actor} & \textbf{Sistema} \\
			\hline
			\begin{enumerate}[noitemsep, nosep]
				\item Ingresa la regla en el campo \IN{regla} desde \UI{configuracion_ui}.
				\item Presiona el botón \BTN{aplicar_regla}.
			\end{enumerate} 
			& 
			\begin{enumerate}[noitemsep, nosep]
				\item[3.] Actualiza la regla.
				\item[4.] Muestra el mensaje \MSG{reglas}. 
			\end{enumerate} \\
			\hline
		\end{tabular} \\
		
		\begin{tabular}{|p{0.45\textwidth}|p{0.45\textwidth}|}
			\hline
			\multicolumn{2}{|c|}{\textbf{Flujo alterno 1: Simulando}} \\ \hline
			\textbf{Actor} & \textbf{Sistema} \\
			\hline
			\begin{enumerate}[noitemsep, nosep]
				\item[2.1.] Presiona el botón \BTN{aplicar_regla} mientras la simulación está activa.
			\end{enumerate} 
			& 
			\begin{enumerate}[noitemsep, nosep]
				\item[3.1.] Pausa la simulación.
				\item[4.1.] Actualiza la regla.
				\item[5.1.] Muestra el mensaje \MSG{reglas}. 
			\end{enumerate} \\
			\hline
		\end{tabular} 
	\end{center}
\end{table} 

\begin{table}[H]
	\begin{center}
		\begin{tabular}{|p{0.2\textwidth}|p{0.7\textwidth}|}
			\hline
			\multicolumn{2}{|c|}{\textbf{Flujo alterno 2: Fuera de rango}} \\ \hline
			\textbf{Actor} & \textbf{Sistema} \\
			\hline
			& 
			\begin{enumerate}[noitemsep, nosep]
				\item[3.3.] La regla no es compatible, se ajusta a la configuración del CA.
				\item[4.3.] Actualiza la regla.
				\item[5.3.] Muestra el mensaje \MSG{reglas}. 
			\end{enumerate} \\
			\hline
		\end{tabular} 
	\end{center}
\end{table}


\subsection{Interacción con el espacio celular}

El usuario puede interactuar con el espacio celular para cambiar valores, colores o vecindades.

\subsubsection{Seleccionar color}

\begin{table}[H]
\begin{center}
	\begin{tabular}{| p{0.2\textwidth} | p{0.7\textwidth} |}
		\hline 
			\textbf{Actor} & Usuario\\ \hline
			\textbf{Descripción} & Se define un color asociado a cada estado.  \\ \hline
			\textbf{Tipo} & Secundario \\ \hline
			\textbf{Entradas} &
			\begin{itemize}[noitemsep, nosep]
			\item Desde mouse \BTN{color}
				\end{itemize}
			\\ \hline
			\textbf{Salidas} & Ninguna  \\ \hline
			\textbf{Precondición} & Ninguna  \\ \hline
			\textbf{Postcondición} &  Ninguna  \\ \hline
			\textbf{Errores} & Ninguno \\ \hline
			\end{tabular}
			\end{center}
			\end{table}
			
		\begin{table}[H]
	\begin{center}
\begin{tabular}{|p{0.45\textwidth}|p{0.45\textwidth}|}
\hline
\multicolumn{2}{|c|}{\textbf{Flujo básico}} \\ \hline
	\textbf{Actor} & \textbf{Sistema} \\
		\hline
			\begin{enumerate}[noitemsep, nosep]
			\item Selecciona un estado en \UI{configuracion_ui}.
			\item Presiona el botón de \BTN{color}.
			\item[4. ] Selecciona un color.
			\end{enumerate} 
				& 
			\begin{enumerate}[noitemsep, nosep]
			\item[3.] Abre un panel de colores. 
			\item[5.] Asigna ese color a ese estado.
				\item[6.] Cambia el color de todas las células con ese estado en el espacio celular.
			\end{enumerate} \\
			\hline
		\end{tabular} \\
	\end{center}
\end{table}

\subsubsection{Cambiar color de estado}

\begin{table}[H]
\begin{center}
	\begin{tabular}{| p{0.2\textwidth} | p{0.7\textwidth} |}
		\hline 
			\textbf{Actor} & Usuario\\ \hline
			\textbf{Descripción} & Se define un color asociado a cada estado.  \\ \hline
			\textbf{Tipo} & Secundario \\ \hline
			\textbf{Entradas} &
			\begin{itemize}[noitemsep, nosep]
			\item Desde mouse \BTN{color}
				\end{itemize}
			\\ \hline
			\textbf{Salidas} & Ninguna  \\ \hline
			\textbf{Precondición} & Ninguna  \\ \hline
			\textbf{Postcondición} &  Ninguna  \\ \hline
			\textbf{Errores} & Ninguno \\ \hline
			\end{tabular}
			\end{center}
			\end{table}
			
		\begin{table}[H]
	\begin{center}
\begin{tabular}{|p{0.45\textwidth}|p{0.45\textwidth}|}
\hline
\multicolumn{2}{|c|}{\textbf{Flujo básico}} \\ \hline
	\textbf{Actor} & \textbf{Sistema} \\
		\hline
			\begin{enumerate}[noitemsep, nosep]
			\item Selecciona un estado en \UI{configuracion_ui}.
			\item Presiona el botón de \BTN{color}.
			\item[4. ] Selecciona un color.
			\end{enumerate} 
				& 
			\begin{enumerate}[noitemsep, nosep]
			\item[3.] Abre un panel de colores. 
			\item[5.] Asigna ese color a ese estado.
				\item[6.] Cambia el color de todas las células con ese estado en el espacio celular.
			\end{enumerate} \\
			\hline
		\end{tabular} \\
	\end{center}
\end{table}

\subsubsection{Cambiar estado de una célula}

\begin{table}[H]
	\begin{center}
		\begin{tabular}{| p{0.2\textwidth} | p{0.7\textwidth} |}
			\hline 
			\textbf{Actor} & Usuario\\ \hline
			\textbf{Descripción} & Se cambia manualmente la configuración del espacio celular, seleccionando el estado de cada célula.  \\ \hline
			\textbf{Tipo} & Primario \\ \hline
			\textbf{Entradas} &
			\begin{itemize}[noitemsep, nosep]
				\item Desde mouse \BTN{celula}
			\end{itemize}
			\\ \hline
			\textbf{Salidas} & Ninguna
			\\ \hline
			\textbf{Precondición} & Ninguna
			\\ \hline
			\textbf{Postcondición} & Ninguna \\ \hline
			\textbf{Errores} & Ninguno
			\\ \hline
		\end{tabular}
	\end{center}
\end{table}

\begin{table}[H]
	\begin{center}
		\begin{tabular}{|p{0.45\textwidth}|p{0.45\textwidth}|}
			\hline
			\multicolumn{2}{|c|}{\textbf{Flujo básico (Siguiente estado)}} \\ \hline
			\textbf{Actor} & \textbf{Sistema} \\
			\hline
			\begin{enumerate}[noitemsep, nosep]
				\item Presiona una célula \BTN{celula} desde \UI{simulacion_ui}.
			\end{enumerate} 
			& 
			\begin{enumerate}[noitemsep, nosep]
				\item[2.] La célula cambia al estado siguiente. 
			\end{enumerate} \\
			\hline
		\end{tabular} \\
	\end{center}
\end{table}

\begin{table}[H]
	\begin{center}
		\begin{tabular}{|p{0.45\textwidth}|p{0.45\textwidth}|}
			\hline
			\multicolumn{2}{|c|}{\textbf{Flujo básico (Selección de estado)}} \\ \hline
			\textbf{Actor} & \textbf{Sistema} \\
			\hline
			\begin{enumerate}[noitemsep, nosep]
				\item Selecciona un estado desde \UI{configuracion_ui}.
				\item Presiona una célula \BTN{celula} desde \UI{simulacion_ui}.
			\end{enumerate} 
			& 
			\begin{enumerate}[noitemsep, nosep]
				\item[3.] La célula cambia al estado presionado en \UI{configuracion_ui}.
			\end{enumerate} \\
			\hline
		\end{tabular} \\
	\end{center}
\end{table}